\section{Motivating Example}

We present an overview of \OurLang.
Modern climate models such as \todo{cite} consist of a large number of interacting modules.
When a data scientist is investigating a new domain, they often begin by analysing
data in an exploratory fashion. This data may be arranged as a collection of records
like in a relational database. Investigating relationships between these complex
datasets can be difficult.

Let us consider the following example: our hypothetical data scientist is 
investigating water-usage in European countries. Each nation has associated with
it a collection of cities and a measure of the total farmland in that nation.
The user may provide a function for computing the total water consumption of a
nation. We picture this
in \figref{water-query}.

\begin{figure}[h]
\tiny
\lstinputlisting[language=Fluid,mathescape]{fig/example/water-query.fld}
\caption{Example Water Usage Query}
\label{fig:water-query}
\end{figure}

With the Galois slicing technique, we can follow a selection on the output of
the program - in this case the integer 570 - back in to the inputs of the program.
We denote a selected by enclosing it in special delimiters like ((so)).